% Options for packages loaded elsewhere
\PassOptionsToPackage{unicode}{hyperref}
\PassOptionsToPackage{hyphens}{url}
\documentclass[
]{article}
\usepackage{xcolor}
\usepackage[margin=1in]{geometry}
\usepackage{amsmath,amssymb}
\setcounter{secnumdepth}{-\maxdimen} % remove section numbering
\usepackage{iftex}
\ifPDFTeX
  \usepackage[T1]{fontenc}
  \usepackage[utf8]{inputenc}
  \usepackage{textcomp} % provide euro and other symbols
\else % if luatex or xetex
  \usepackage{unicode-math} % this also loads fontspec
  \defaultfontfeatures{Scale=MatchLowercase}
  \defaultfontfeatures[\rmfamily]{Ligatures=TeX,Scale=1}
\fi
\usepackage{lmodern}
\ifPDFTeX\else
  % xetex/luatex font selection
\fi
% Use upquote if available, for straight quotes in verbatim environments
\IfFileExists{upquote.sty}{\usepackage{upquote}}{}
\IfFileExists{microtype.sty}{% use microtype if available
  \usepackage[]{microtype}
  \UseMicrotypeSet[protrusion]{basicmath} % disable protrusion for tt fonts
}{}
\makeatletter
\@ifundefined{KOMAClassName}{% if non-KOMA class
  \IfFileExists{parskip.sty}{%
    \usepackage{parskip}
  }{% else
    \setlength{\parindent}{0pt}
    \setlength{\parskip}{6pt plus 2pt minus 1pt}}
}{% if KOMA class
  \KOMAoptions{parskip=half}}
\makeatother
\usepackage{color}
\usepackage{fancyvrb}
\newcommand{\VerbBar}{|}
\newcommand{\VERB}{\Verb[commandchars=\\\{\}]}
\DefineVerbatimEnvironment{Highlighting}{Verbatim}{commandchars=\\\{\}}
% Add ',fontsize=\small' for more characters per line
\usepackage{framed}
\definecolor{shadecolor}{RGB}{248,248,248}
\newenvironment{Shaded}{\begin{snugshade}}{\end{snugshade}}
\newcommand{\AlertTok}[1]{\textcolor[rgb]{0.94,0.16,0.16}{#1}}
\newcommand{\AnnotationTok}[1]{\textcolor[rgb]{0.56,0.35,0.01}{\textbf{\textit{#1}}}}
\newcommand{\AttributeTok}[1]{\textcolor[rgb]{0.13,0.29,0.53}{#1}}
\newcommand{\BaseNTok}[1]{\textcolor[rgb]{0.00,0.00,0.81}{#1}}
\newcommand{\BuiltInTok}[1]{#1}
\newcommand{\CharTok}[1]{\textcolor[rgb]{0.31,0.60,0.02}{#1}}
\newcommand{\CommentTok}[1]{\textcolor[rgb]{0.56,0.35,0.01}{\textit{#1}}}
\newcommand{\CommentVarTok}[1]{\textcolor[rgb]{0.56,0.35,0.01}{\textbf{\textit{#1}}}}
\newcommand{\ConstantTok}[1]{\textcolor[rgb]{0.56,0.35,0.01}{#1}}
\newcommand{\ControlFlowTok}[1]{\textcolor[rgb]{0.13,0.29,0.53}{\textbf{#1}}}
\newcommand{\DataTypeTok}[1]{\textcolor[rgb]{0.13,0.29,0.53}{#1}}
\newcommand{\DecValTok}[1]{\textcolor[rgb]{0.00,0.00,0.81}{#1}}
\newcommand{\DocumentationTok}[1]{\textcolor[rgb]{0.56,0.35,0.01}{\textbf{\textit{#1}}}}
\newcommand{\ErrorTok}[1]{\textcolor[rgb]{0.64,0.00,0.00}{\textbf{#1}}}
\newcommand{\ExtensionTok}[1]{#1}
\newcommand{\FloatTok}[1]{\textcolor[rgb]{0.00,0.00,0.81}{#1}}
\newcommand{\FunctionTok}[1]{\textcolor[rgb]{0.13,0.29,0.53}{\textbf{#1}}}
\newcommand{\ImportTok}[1]{#1}
\newcommand{\InformationTok}[1]{\textcolor[rgb]{0.56,0.35,0.01}{\textbf{\textit{#1}}}}
\newcommand{\KeywordTok}[1]{\textcolor[rgb]{0.13,0.29,0.53}{\textbf{#1}}}
\newcommand{\NormalTok}[1]{#1}
\newcommand{\OperatorTok}[1]{\textcolor[rgb]{0.81,0.36,0.00}{\textbf{#1}}}
\newcommand{\OtherTok}[1]{\textcolor[rgb]{0.56,0.35,0.01}{#1}}
\newcommand{\PreprocessorTok}[1]{\textcolor[rgb]{0.56,0.35,0.01}{\textit{#1}}}
\newcommand{\RegionMarkerTok}[1]{#1}
\newcommand{\SpecialCharTok}[1]{\textcolor[rgb]{0.81,0.36,0.00}{\textbf{#1}}}
\newcommand{\SpecialStringTok}[1]{\textcolor[rgb]{0.31,0.60,0.02}{#1}}
\newcommand{\StringTok}[1]{\textcolor[rgb]{0.31,0.60,0.02}{#1}}
\newcommand{\VariableTok}[1]{\textcolor[rgb]{0.00,0.00,0.00}{#1}}
\newcommand{\VerbatimStringTok}[1]{\textcolor[rgb]{0.31,0.60,0.02}{#1}}
\newcommand{\WarningTok}[1]{\textcolor[rgb]{0.56,0.35,0.01}{\textbf{\textit{#1}}}}
\usepackage{longtable,booktabs,array}
\newcounter{none} % for unnumbered tables
\usepackage{calc} % for calculating minipage widths
% Correct order of tables after \paragraph or \subparagraph
\usepackage{etoolbox}
\makeatletter
\patchcmd\longtable{\par}{\if@noskipsec\mbox{}\fi\par}{}{}
\makeatother
% Allow footnotes in longtable head/foot
\IfFileExists{footnotehyper.sty}{\usepackage{footnotehyper}}{\usepackage{footnote}}
\makesavenoteenv{longtable}
\usepackage{graphicx}
\makeatletter
\newsavebox\pandoc@box
\newcommand*\pandocbounded[1]{% scales image to fit in text height/width
  \sbox\pandoc@box{#1}%
  \Gscale@div\@tempa{\textheight}{\dimexpr\ht\pandoc@box+\dp\pandoc@box\relax}%
  \Gscale@div\@tempb{\linewidth}{\wd\pandoc@box}%
  \ifdim\@tempb\p@<\@tempa\p@\let\@tempa\@tempb\fi% select the smaller of both
  \ifdim\@tempa\p@<\p@\scalebox{\@tempa}{\usebox\pandoc@box}%
  \else\usebox{\pandoc@box}%
  \fi%
}
% Set default figure placement to htbp
\def\fps@figure{htbp}
\makeatother
\usepackage{svg}
\setlength{\emergencystretch}{3em} % prevent overfull lines
\providecommand{\tightlist}{%
  \setlength{\itemsep}{0pt}\setlength{\parskip}{0pt}}
\usepackage{bookmark}
\IfFileExists{xurl.sty}{\usepackage{xurl}}{} % add URL line breaks if available
\urlstyle{same}
% fallback for those not using the hyperref driver hyperxmp:
\makeatletter
\@ifundefined{xmpquote}{\newcommand{\xmpquote}[1]{#1}}{}
\makeatother
\hypersetup{
  hidelinks,
  pdfcreator={LaTeX via pandoc}}

\author{}
\date{\vspace{-2.5em}}

\begin{document}

\section{modelterms: Coerce, Inspect, Filter, and Match Model Terms in
R}\label{modelterms-coerce-inspect-filter-and-match-model-terms-in-r}

\href{https://lifecycle.r-lib.org/articles/stages.html\#experimental}{\pandocbounded{\includesvg[keepaspectratio]{https://img.shields.io/badge/lifecycle-experimental-orange.svg}}}

\textbf{modelterms} is a lightweight, zero-dependency R package
providing a unified interface for manipulating \texttt{stats::terms}
objects, formulas, fitted models, and character vectors of term labels.

It is designed both as a standalone utility for statistical modeling
workflows and as a helper module for other packages (including future
versions of \href{https://github.com/ryo-asashiba/midr}{\textbf{midr}}).

\subsection{\texorpdfstring{Why
\textbf{modelterms}?}{Why modelterms?}}\label{why-modelterms}

In R, model terms serve as the structural blueprint for predictive
model, a role of paramount importance for additive models built with
\texttt{stats::lm()}, \texttt{stats::glm()}, or GAMs. However, the term
``terms'' has historically suffered from ambiguity. For instance, it
appears as the \texttt{terms} argument in \texttt{stats::predict.lm()}
which accepts character vectors of additive components, whereas the
\texttt{"terms"} objects stored inside fitted models are formula packed
with complex attributes, not character vectors.

While Base R's \texttt{stats::terms()} and \texttt{stats::drop.terms()}
provide powerful formula parsing, working with them programmatically
often involves sharp edges:

\begin{enumerate}
\def\labelenumi{\arabic{enumi}.}
\tightlist
\item
  \textbf{Inconsistent Coercion}: \texttt{stats::terms()} cannot
  directly coerce a character vector of term labels like
  \texttt{c("x1",\ "x2")}, and \texttt{stats::reformulate(character(0))}
  throws an error.
\item
  \textbf{Unwanted Reordering \& Renaming}: Even with
  \texttt{keep.order\ =\ TRUE}, \texttt{stats::terms()} silently
  rewrites interaction labels based on variable appearance order (e.g.,
  turning \texttt{c("a",\ "b:a")} into \texttt{c("a",\ "a:b")}).
\item
  \textbf{Fragile Subsetting}: \texttt{stats::drop.terms()} requires
  numeric indices, loses the response variable by default, and crashes
  if all terms are dropped.
\item
  \textbf{Ambiguous ``Variables''}: Base R lacks a unified way to
  distinguish between \emph{term labels} (term names in additive models,
  e.g., \texttt{"x1:log(x2)"}), \emph{model variables} (transformed
  variable names in formula, e.g., \texttt{"log(x2)"}), and underlying
  \emph{raw variables} (column names in data.frame, e.g.,
  \texttt{"x2"}).
\end{enumerate}

\textbf{modelterms} resolves all of these issues with a clean,
pipe-friendly functional API.

\subsection{Installation}\label{installation}

You can install the development version from GitHub:

\begin{Shaded}
\begin{Highlighting}[]
\CommentTok{\# install.packages("pak")}
\CommentTok{\# pak::pak("ryo{-}asashiba/modelterms")}
\end{Highlighting}
\end{Shaded}

\subsection{Core Concepts: Three Levels of Model
Components}\label{core-concepts-three-levels-of-model-components}

Consider the formula
\texttt{y\ \textasciitilde{}\ x1\ +\ log(x2)\ +\ x1:log(x2)\ +\ I(x1\ *\ x2)}.
This package clearly distinguishes three structural levels:

{\def\LTcaptype{none} % do not increment counter
\begin{longtable}[]{@{}
  >{\raggedright\arraybackslash}p{(\linewidth - 6\tabcolsep) * \real{0.2500}}
  >{\raggedright\arraybackslash}p{(\linewidth - 6\tabcolsep) * \real{0.2500}}
  >{\raggedright\arraybackslash}p{(\linewidth - 6\tabcolsep) * \real{0.2500}}
  >{\raggedright\arraybackslash}p{(\linewidth - 6\tabcolsep) * \real{0.2500}}@{}}
\toprule\noalign{}
\begin{minipage}[b]{\linewidth}\raggedright
Level
\end{minipage} & \begin{minipage}[b]{\linewidth}\raggedright
Concept
\end{minipage} & \begin{minipage}[b]{\linewidth}\raggedright
Function
\end{minipage} & \begin{minipage}[b]{\linewidth}\raggedright
Output for Example Formula
\end{minipage} \\
\midrule\noalign{}
\endhead
\bottomrule\noalign{}
\endlastfoot
\textbf{Term Labels} & Individual additive components (main effects \&
interactions) & \texttt{get\_labels(x)} &
\texttt{c("x1",\ "log(x2)",\ "I(x1\ *\ x2)",\ "x1:log(x2)")} \\
\textbf{Model Variables} & Distinct predictor expressions (rows of the
\texttt{factors} matrix) & \texttt{get\_variables(x,\ raw\ =\ FALSE)} &
\texttt{c("x1",\ "log(x2)",\ "I(x1\ *\ x2)")} \\
\textbf{Raw Variables} & Underlying column names required in a
\texttt{data.frame} & \texttt{get\_variables(x,\ raw\ =\ TRUE)} &
\texttt{c("x1",\ "x2")} \\
\end{longtable}
}

Setting \texttt{each.term\ =\ TRUE} in \texttt{get\_variables()} returns
a named list mapping each term label to its constituent variables.

\subsection{Usage Examples}\label{usage-examples}

\subsubsection{\texorpdfstring{1. Flexible and Safe Coercion with
\texttt{as\_terms()}}{1. Flexible and Safe Coercion with as\_terms()}}\label{flexible-and-safe-coercion-with-as_terms}

\texttt{as\_terms()} accepts \texttt{terms} objects, formulas, fitted
models (\texttt{lm}, \texttt{glm}, etc.), or character vectors.

\begin{itemize}
\tightlist
\item
  \textbf{Strict Term Label Vectors
  (\texttt{length\ \textgreater{}=\ 2})}: Every element must be a valid
  single term. Input order and exact interaction notation (such as
  \texttt{"b:a"}) are strictly preserved.
\item
  \textbf{Formula Shorthand (\texttt{length\ ==\ 1})}: A single string
  containing formula operators (e.g., \texttt{".\^{}2"} or
  \texttt{"a\ *\ b\ +\ c"}) is expanded as a formula right-hand side.
\end{itemize}

\begin{Shaded}
\begin{Highlighting}[]
\FunctionTok{library}\NormalTok{(modelterms)}

\CommentTok{\# Preserves exact order and "b:a" notation for term label vectors}
\NormalTok{t1 }\OtherTok{\textless{}{-}} \FunctionTok{as\_terms}\NormalTok{(}\FunctionTok{c}\NormalTok{(}\StringTok{"b:a"}\NormalTok{, }\StringTok{"a"}\NormalTok{, }\StringTok{"1"}\NormalTok{))}
\FunctionTok{get\_labels}\NormalTok{(t1)}
\end{Highlighting}
\end{Shaded}

\begin{verbatim}
## [1] "b:a" "a"
\end{verbatim}

\begin{Shaded}
\begin{Highlighting}[]
\FunctionTok{has\_intercept}\NormalTok{(t1)}
\end{Highlighting}
\end{Shaded}

\begin{verbatim}
## [1] TRUE
\end{verbatim}

\begin{Shaded}
\begin{Highlighting}[]
\CommentTok{\# Expands length{-}1 formula shorthands (including "." with data)}
\NormalTok{t2 }\OtherTok{\textless{}{-}} \FunctionTok{as\_terms}\NormalTok{(}\StringTok{".\^{}2"}\NormalTok{, }\AttributeTok{y =} \StringTok{"Ozone"}\NormalTok{, }\AttributeTok{data =}\NormalTok{ airquality[, }\DecValTok{1}\SpecialCharTok{:}\DecValTok{4}\NormalTok{])}
\FunctionTok{get\_labels}\NormalTok{(t2)}
\end{Highlighting}
\end{Shaded}

\begin{verbatim}
## [1] "Solar.R"      "Wind"         "Temp"         "Solar.R:Wind" "Solar.R:Temp"
## [6] "Wind:Temp"
\end{verbatim}

\subsubsection{2. Filtering Terms and
Plotting}\label{filtering-terms-and-plotting}

\texttt{filter\_terms()} subsets a ``terms'' object by interaction
order, required variables, or excluded variables. Unlike
\texttt{stats::drop.terms()}, it safely returns an intercept-only (or
empty) ``terms'' object even when all terms are filtered out.

\begin{Shaded}
\begin{Highlighting}[]
\NormalTok{fml }\OtherTok{\textless{}{-}} \FunctionTok{log}\NormalTok{(Ozone) }\SpecialCharTok{\textasciitilde{}}\NormalTok{ Solar.R }\SpecialCharTok{+}\NormalTok{ Wind }\SpecialCharTok{+}\NormalTok{ Temp }\SpecialCharTok{+}\NormalTok{ Wind}\SpecialCharTok{:}\NormalTok{Temp }\SpecialCharTok{+}\NormalTok{ Solar.R}\SpecialCharTok{:}\NormalTok{Wind}

\CommentTok{\# Extract main effects only as a character vector}
\NormalTok{fml }\SpecialCharTok{|\textgreater{}} \FunctionTok{filter\_terms}\NormalTok{(}\AttributeTok{order =} \DecValTok{1}\NormalTok{) }\SpecialCharTok{|\textgreater{}} \FunctionTok{get\_labels}\NormalTok{()}
\end{Highlighting}
\end{Shaded}

\begin{verbatim}
## [1] "Solar.R" "Wind"    "Temp"
\end{verbatim}

\begin{Shaded}
\begin{Highlighting}[]
\CommentTok{\#\textgreater{} [1] "Solar.R" "Wind"    "Temp"}

\CommentTok{\# Keep terms involving "Wind", but exclude any term involving "Solar.R"}
\NormalTok{sub\_terms }\OtherTok{\textless{}{-}} \FunctionTok{filter\_terms}\NormalTok{(fml, }\AttributeTok{require =} \StringTok{"Wind"}\NormalTok{, }\AttributeTok{remove =} \StringTok{"Solar.R"}\NormalTok{)}
\FunctionTok{get\_labels}\NormalTok{(sub\_terms)}
\end{Highlighting}
\end{Shaded}

\begin{verbatim}
## [1] "Wind"      "Wind:Temp"
\end{verbatim}

\begin{Shaded}
\begin{Highlighting}[]
\CommentTok{\#\textgreater{} [1] "Wind"      "Wind:Temp"}

\CommentTok{\# Directly use the filtered terms object to fit a submodel}
\NormalTok{lm\_sub }\OtherTok{\textless{}{-}} \FunctionTok{lm}\NormalTok{(sub\_terms, }\AttributeTok{data =}\NormalTok{ airquality)}
\end{Highlighting}
\end{Shaded}

\subsubsection{3. Order-Invariant Term
Matching}\label{order-invariant-term-matching}

When users specify interaction terms in plotting or evaluation
functions, they may write \texttt{"Temp:Wind"} instead of
\texttt{"Wind:Temp"}. \texttt{match\_labels()} normalizes and matches
terms regardless of variable order while rejecting multi-term
expressions like \texttt{"Wind\ *\ Temp"}.

\begin{Shaded}
\begin{Highlighting}[]
\NormalTok{fml }\OtherTok{\textless{}{-}}\NormalTok{ y }\SpecialCharTok{\textasciitilde{}}\NormalTok{ x1 }\SpecialCharTok{+} \FunctionTok{log}\NormalTok{(x2) }\SpecialCharTok{+}\NormalTok{ x1}\SpecialCharTok{:}\FunctionTok{log}\NormalTok{(x2) }\SpecialCharTok{+} \FunctionTok{I}\NormalTok{(}\DecValTok{1}\SpecialCharTok{:}\DecValTok{100}\NormalTok{)}\SpecialCharTok{:}\NormalTok{x1}

\CommentTok{\# Matches reversed interaction notation to the canonical model labels}
\FunctionTok{match\_labels}\NormalTok{(}\FunctionTok{c}\NormalTok{(}\StringTok{"log(x2):x1"}\NormalTok{, }\StringTok{"x1:I(1:100)"}\NormalTok{, }\StringTok{"x1"}\NormalTok{), fml)}
\end{Highlighting}
\end{Shaded}

\begin{verbatim}
## [1] "x1:log(x2)"  "x1:I(1:100)" "x1"
\end{verbatim}

\begin{Shaded}
\begin{Highlighting}[]
\CommentTok{\#\textgreater{} [1] "x1:log(x2)"  "I(1:100):x1" "x1"}

\CommentTok{\# Safely handles unmatched or invalid expressions ("pass", "omit", or "fail")}
\FunctionTok{match\_labels}\NormalTok{(}\FunctionTok{c}\NormalTok{(}\StringTok{"x1"}\NormalTok{, }\StringTok{"x1 * x2"}\NormalTok{, }\StringTok{"unknown"}\NormalTok{), fml, }\AttributeTok{nomatch =} \StringTok{"omit"}\NormalTok{)}
\end{Highlighting}
\end{Shaded}

\begin{verbatim}
## [1] "x1"
\end{verbatim}

\begin{Shaded}
\begin{Highlighting}[]
\CommentTok{\#\textgreater{} [1] "x1"}
\end{Highlighting}
\end{Shaded}

\subsection{Function Reference
Summary}\label{function-reference-summary}

\begin{itemize}
\tightlist
\item
  Coercion: \texttt{as\_terms()}
\item
  Filtering: \texttt{filter\_terms()}
\item
  Extraction: \texttt{get\_labels()}, \texttt{get\_variables()},
  \texttt{get\_response()}, \texttt{get\_order()}
\item
  Inspection \& Predicates: \texttt{has\_intercept()},
  \texttt{has\_response()}, \texttt{is\_single\_term()},
  \texttt{is\_fully\_crossed()}
\item
  Matching: \texttt{make\_signature()}, \texttt{match\_labels()}
\end{itemize}

\end{document}
